\subsubsection{Prompt Engineering} \label{sec:llm_prompt_engineering}

Prompt Engineering ist eine entscheidende Methodik in der Arbeit mit \acp{llm}, die jede Person bewusst oder unterbewusst bei jeder Anfrage verwendet. Ein Prompt (\textit{dt. Eingabeaufforderung}) ist eine aufgabenspezifische Textanweisung, die dem Modell zur Inferenz übergeben wird, um die Ausgabe zu steuern, ohne die Modellparameter zu verändern. Im Kern des Prompt Engineering steht das Bestreben, durch methodisch geformte Texteingaben die Ausgaben künstlicher Intelligenz zu optimieren. \cite{prompt-survey}

Zu den bekanntesten Techniken des Prompt Engineering gehören \cite{prompt-survey}:
\begin{itemize}
    \item \textbf{Zero-Shot Prompting:} Bei dieser Technik wird das Modell mit einer Aufgabe konfrontiert, ohne dass es vorherige Beispiele oder Anweisungen erhält. Es wird erwartet, dass das Modell die Aufgabe aufgrund des vorhandenen Wissens und des gegebenen Textes versteht und ausführt.
    Beispiel Prompt: "\texttt{Schreibe einen Text über die Vor- und Nachteile von \acp{llm}.}"
    \item \textbf{Few-Shot Prompting:} Hierbei werden dem Modell einige Beispiele der Aufgabe gegeben, um ihm zu helfen, die Struktur und den Kontext der gewünschten Ausgabe zu verstehen. Diese Technik ist besonders nützlich, wenn das Modell Schwierigkeiten hat, die Aufgabe allein zu bewältigen oder aus vorgegebenen Auswahlmöglichkeiten auswählen muss. Allerdings können durch die gegebenen Beispiele häufige Wörter bevorzugt werden oder die Kreativität sinken.
    Beispiel Prompt: "\texttt{Beispiel 1: Frage: Was ist die Hauptstadt von Frankreich? Antwort: Paris. Beispiel 2: Frage: Was ist die Hauptstadt von Italien? Antwort: Rom. Aufgabe: Was ist die Hauptstadt von Spanien? Antwort:}"
    \item \textbf{Chain-of-Thought Prompting:} Diese Technik ermutigt das Modell, seine Schlussfolgerungen und Überlegungen während des Problemlösungsprozesses zu erläutern. Diese Methode generiert eine Kette von Gedanken, also Argumentationsschritte, um so Probleme schrittweise zu lösen.
    Beispiel Prompt: "\texttt{Schritt 1: Erkläre was \acp{llm} sind. Schritt 2: Erkläre welche Einsatzgebiete sie haben. Schritt 3: Erkläre wie sie funktionieren. Basierend auf den Schritten, erkläre was die Vor- und Nachteile sind.}"
    \item \textbf{Role Prompting:} Bei dieser Technik wird das Modell in eine bestimmte Rolle versetzt, um die Ausgabe zu steuern. Dies kann durch die Verwendung von Expertenrollen geschehen, die das Modell dazu bringen, aus der Perspektive einer bestimmten Person oder Entität zu antworten. \cite{prompt-role}
    Beispiel Prompt: "\texttt{Du bist ein guter Grundschullehrer und kannst Probleme einfach runterbrechen. Erkläre einem Grundschüler, was \acp{llm} sind und wie sie funktionieren.}"
    \item \textbf{Self-Consistency Prompting:} Diese Technik ermutigt das Modell, mehrere Antworten auf eine Frage zu generieren und dann die beste Antwort auszuwählen. Für komplexe Aufgaben mit mehreren validen Lösungswegen, werden diese verglichen und die schlüssigste Antwort gewählt. 
    Beispiel Prompt: "\texttt{Erkläre was \acp{llm} sind. Generiere 5 verschiedene Erklärungen und wähle die schlüssigste und klarste aus.}"
\end{itemize}

In dem Prompting Framework "TACOMORE" (Task, Context, Model, Reproducibility) gehört neben dem Definieren von einer oder mehreren Aufgaben mittels oben genannten Techniken auch das Bereitstellen von passenden Kontextinformationen, die Auswahl des richtigen Modells und die Definition des gewünschten Ausgabeformats, um Reproduzierbarkeit sicherzustellen. Wegen der Natur von \acp{llm} kann nicht bei jeder Ausführung desselben Prompts dieselbe Ausgabe garantiert werden. Deswegen ist es wichtig, dass man das Ausgabeformat festlegt. Maschinenlesbare Formate wie JSON oder XML sind passend, da der Datentyp in einem Schlüssel-Wert-Verfahren festgelegt werden kann. \cite{prompt-tacomore}

Durch die Kombination der oben genannten Techniken und Frameworks kann die Leistung und die Genauigkeit erheblich verbessert werden. Es ist wichtig, dass die Prompt Engineering Techniken an die spezifischen Anforderungen und Ziele der Anwendung angepasst werden, um die besten Ergebnisse zu erzielen. Besonders bei Anwendungsfällen, wo die Nutzenden die Ausgabe nicht direkt sehen und keine Möglichkeit haben, weitere Prompts nachzusenden, um die Ausgabe zu verbessern, ist Prompt Engineering entscheidend. Das Ausgabeformat und die Qualität müssen auf den ersten Versuch stimmen. Sonst kann es zu Unzufriedenheit oder sogar Fehlern kommen. So ein Anwendungsfall ist unter anderem die Integration von KI in ein Umfrage-Tool.